% =====================================================================
% Capítulo 5 -- Resultados (versão para a qualificação: resultados esperados, plano, cronograma e riscos)
%
% Nesta etapa NÃO há resultados de modelagem. O capítulo registra o que já se
% sabe (Capítulo 4), as decisões em aberto e o plano de modelagem e validação.
% Na defesa, o título volta a ser "Resultados" e a seção 5.2 é substituída
% pelos resultados obtidos, e as seções 5.5 (cronograma) e 5.6 (riscos) saem.
% \pendente = decisão da autora ou do orientador.
% Usa os rótulos chp:ead (cap. 4) e chp:fundamentacao (cap. 2) do arquivo principal.
% =====================================================================

\section{Situação desta etapa}
\label{sec:res-situacao}

Até a data da qualificação, a modelagem preditiva ainda não foi executada, e por isso este capítulo não apresenta resultados de modelos. O que existe é a base de dados integrada e a análise exploratória descritiva (Capítulo~\ref{chp:ead}), que caracterizou o Atlas Digital de Desastres, o ERA5 e o MapBiomas nos 35 municípios costeiros de Santa Catarina entre 1991 e 2025. Este capítulo reúne o que essa análise permite antecipar, as decisões que ela deixou em aberto e o plano de modelagem e de validação que será seguido a partir do cronograma (Tabela~\ref{tab:cronograma}, Seção~\ref{sec:res-cronograma}) e dos riscos associados (Seção~\ref{sec:res-riscos}).

\section{O que a análise exploratória antecipa}
\label{sec:res-antecipa}

A análise exploratória não testa as hipóteses de pesquisa, mas indica como elas devem ser testadas. Três achados orientam o plano.

O primeiro é que a variável resposta muda com a versão escolhida. O total de eventos hidrológicos cresce cerca de três vezes depois de 2013, mas, sem a tipologia ``Chuvas Intensas'', criada nesse ano, a média anual praticamente não muda (razão de 1,13). Assim, qualquer resultado de previsão deve dizer de qual versão da resposta trata, e comparar versões é parte do desenho.

O segundo é que a chuva mensal da reanálise não mostra tendência no período, enquanto os registros sobem. Isso indica que o aumento dos registros não acompanha, nessa escala, o aumento da chuva, e que parte dele pode vir de mudanças na forma de registrar os desastres. Por isso, a avaliação da hipótese de que a inteligência artificial melhora a previsão precisa separar o período anterior e o posterior a 2013.

O terceiro é que as variáveis de uso da terra variam quase só entre municípios e se misturam com a latitude. Elas descrevem diferenças entre lugares, e não mudanças no tempo, e precisam entrar no modelo como características do município.

\section{Decisões em aberto}
\label{sec:res-decisoes}

O Capítulo~\ref{chp:ead} deixou decisões que mudam os números do que será modelado: a versão da variável resposta (total, sem Chuvas Intensas ou ambas); o corte do período (2013, primeiro ano de Chuvas Intensas, ou 2012); a conversão de unidades da precipitação do ERA5; a escolha do vento como controle negativo ou de outra variável que não covarie com a chuva; e a ausência do MapBiomas em Camboriú e Paulo Lopes, que parece falha de junção \pendente{registrar a decisão tomada em cada item antes da qualificação}.

A segunda hipótese de pesquisa, de que variáveis socioeconômicas ajudam a identificar áreas críticas, depende de dados que ainda não estão na base: hoje ela tem chuva, vento, uso da terra e áreas de risco, e não tem indicadores socioeconômicos \pendente{definir a fonte dos indicadores socioeconômicos e o índice de vulnerabilidade, ou ajustar a hipótese}.

\section{Plano de modelagem e validação}
\label{sec:res-plano}

A modelagem seguirá três cuidados que decorrem da análise exploratória.

\paragraph*{Modelos de referência primeiro.} Antes de árvores de decisão, redes neurais ou LSTM (Capítulo~\ref{chp:fundamentacao}), serão ajustados modelos simples de referência, como a média histórica do município. Um modelo mais complexo só é considerado melhor se superar esses modelos, o que dá sentido à primeira hipótese.

\paragraph*{Validação que respeita o tempo e o espaço.} Como os registros mudam de forma em 2013 e municípios vizinhos têm valores parecidos (índice de Moran de 0,37 na densidade de eventos), a divisão aleatória entre treino e teste superestimaria o desempenho. A validação usará períodos posteriores como teste e, em separado, municípios deixados de fora \pendente{confirmar os blocos de tempo e de município com o orientador}.

\paragraph*{Métricas adequadas a poucos eventos.} Em 62,8\% dos município-anos não há registro de evento, e a série é de apenas 35 anos. Por isso as métricas devem tratar eventos raros e contagens, e não apenas o erro médio \pendente{definir as métricas; por exemplo, o erro absoluto médio da contagem e a área sob a curva de precisão e revocação da ocorrência de evento}.

\section{Cronograma}
\label{sec:res-cronograma}

A Tabela~\ref{tab:cronograma} apresenta o cronograma restante até a defesa. A coleta e organização dos dados e a análise exploratória (Capítulo~\ref{chp:ead}) já foram concluídas em M19. A implementação dos modelos ocupa os meses M20 e M21, a validação e os testes ocupam M21 e M22, e a escrita final e a defesa ocupam M22 a M24. As decisões em aberto da Seção~\ref{sec:res-decisoes} precisam ser fechadas no início de M20, porque definem a variável resposta e o corte do período usados na modelagem.

\begin{table}[htbp]
\centering
\caption{Cronograma restante para a defesa}
\label{tab:cronograma}
\begin{tabular}{lcccccc}
\hline
Atividade & M19 & M20 & M21 & M22 & M23 & M24 \\
\hline
Levantamento bibliográfico & X & X & X & X & X & \\
Coleta e organização dos dados & X & & & & & \\
Análise exploratória dos dados & X & & & & & \\
Decisões em aberto sobre a base (Seção~\ref{sec:res-decisoes}) & & X & & & & \\
Implementação do modelo & & X & X & & & \\
Validação e testes & & & X & X & & \\
Escrita da dissertação & & & & X & X & \\
Defesa & & & & & & X \\
\hline
\end{tabular}
\par\small Fonte: Elaborado pela autora.
\end{table}

\section{Análise de riscos}
\label{sec:res-riscos}

A Tabela~\ref{tab:riscos} reúne os riscos do projeto. Os cinco primeiros vêm da proposta original; os quatro últimos surgiram da análise exploratória e já têm efeito sobre o plano de modelagem \pendente{revisar probabilidade e impacto de cada risco com o orientador}.

\begin{table}[htbp]
\centering
\caption{Análise de riscos do projeto}
\label{tab:riscos}
\small
\begin{tabular}{p{3.3cm}cc p{3.8cm}p{4.6cm}}
\hline
Risco potencial & Prob. & Impacto & Gatilho & Plano de contingência \\
\hline
Falta de dados históricos completos & Média & Alta & Dados incompletos ou inconsistentes & Buscar fontes alternativas, interpolação e limpeza de dados \\
Baixa performance dos modelos de inteligência artificial & Média & Média & Alto erro ou baixa precisão & Testar algoritmos alternativos, ajustar hiperparâmetros \\
Dificuldade na integração de dados heterogêneos & Alta & Média & Inconsistências entre fontes & Uso de ferramentas de ETL e padronização \\
Mudanças na disponibilidade de bases públicas & Baixa & Média & Remoção de acesso a APIs ou sites & Armazenar cópias locais, antecipar a coleta de dados \\
Atraso no cronograma devido à complexidade & Média & Alta & Atividades técnicas levarem mais tempo & Reorganizar etapas e priorizar entregas parciais \\
Mudança no registro de desastres ao longo do tempo (Chuvas Intensas a partir de 2013) & Alta & Alta & Total de eventos cresce cerca de três vezes após 2013 sem mudança na chuva & Modelar versões da resposta separadamente, validar por período e relatar a versão usada \\
Poucos eventos por município-ano (62,8\% sem registro) & Alta & Média & Modelos preveem sempre zero ou superam o baseline só por acaso & Métricas para eventos raros, modelos de contagem e comparação com modelos de referência \\
Ausência de indicadores socioeconômicos na base & Alta & Alta & Segunda hipótese sem dados para ser testada & Definir a fonte e o índice de vulnerabilidade em M20, ou ajustar a hipótese \\
Falhas de unidade ou de junção nas bases (ERA5, MapBiomas) & Média & Média & Valores fora da faixa esperada, municípios sem dados & Conferir a conversão do ERA5 e a junção do MapBiomas antes de modelar \\
\hline
\end{tabular}
\par\small Fonte: Elaborado pela autora.
\end{table}

Os resultados obtidos nessas etapas, e não este plano, serão apresentados neste capítulo na versão da defesa.